\scanpage{80}

% L01 -- top, hypothesis
On suppose qu'on est sur \uline{corps $k$}.

% L02 -- upper third, theorem
\uline{Théorème} Soit $Z\subset X$ sous-espace projectif. \deleted{\unsure{Soit}} Pour que $Z$ soit invariant par $G$ i.e. par les $\tau_i$ \inserted{$(0\le i\le n)$}, il f et s. qu'on soit dans l'un des cas \unsure{suivant|suivants} :

% L03 -- upper third, a
a) $Z$ est égal à l'un des ss-espaces $\varnothing,X,\{c\},C$.

% L04 -- upper third, b
b) $\exists\ 0\le i\le n-1$ tel que $\beta_i\ (=2+\alpha_i)=0$, et on a

% L05 -- upper middle, two alternatives
\[
Z=\underbrace{H_0\cap H_1\cap\cdots\cap H_i}_{Z'_i\ (\text{de dim }n-i)}
\quad\text{ou}\quad
Z=\deleted{\unsure{L}}\underbrace{\mathrm{Lin}(h_{i+1},\ldots,h_n)}_
 {\substack{Z_i^{\deleted{\unsure{\prime|*}}}\\(\dim n-i-1)}}
\]

% L06 -- left beside alternatives
% slanted
\margin{NB $\beta_i=0\ \Longleftrightarrow\ \underbrace{Z_i\subset Z'_i}_{\text{codim }1}$}

% L07 -- left of b, isolated
\margin{\unsure{b}}

% L08 -- centre, convention
\[
\begin{aligned}
(\text{\uline{NB} On posera}\quad Z_n&=\underbrace{H_0\cap\cdots\cap H_n}_{\{c\}}
 &\text{et}\quad Z'_n&=\varnothing\ \deleted{\unsure{H_n}}.\\
 Z_{-1}&=X &Z'_{-1}&=C
\end{aligned}
\]

% L09 -- above two chains
On a (sans hyp. sur \unsure{aucun|aucune} des $\beta_i$)

% L10 -- middle, two horizontal chains with vertical inclusions
\diagram{\resizebox{\linewidth}{!}{\begin{tikzcd}[ampersand replacement=\&,column sep=2pt,row sep=1pt]
 Z_n^{\deleted{\prime}}=\varnothing \& \subset \& Z_{n-1}^{\deleted{\prime}}=\{h_n\} \& \subset \& Z_{n-2}^{\deleted{\prime}}=\unsure{\mathrm L|\mathrm{Lin}}(h_n,h_{n-1}) \& \cdots \& \subset \& Z_0^{\deleted{\prime}}=\mathrm L(h_n,\ldots,h_1) \& \subset \& Z_{-1}^{\deleted{\prime}}=C\\
 \cap \& {} \& \cap \& {} \& \cap \& {} \& {} \& \cap \& {} \& \cap\\
 Z'_n=\{c\} \& \subset \& Z'_{n-1}=H_0\cap\cdots\cap H_{n-1} \& \subset \& Z'_{n-2}=H_0\cap\cdots\cap H_{n-2} \& \cdots \& \subset \& Z'_0=H_0 \& \subset \& Z'_{-1}=X
\end{tikzcd}}}

% L11 -- below chains
Toutes les inclusions écrites sont strictes, de codim 1.

% L12 -- lower third, dimensions
\deleted{Tous} les $2(n+2)$ espaces $Z_i,Z'_i$ $( -\unsure{0|1}\le i\le n)$ sont\inserted{-ils} distincts ? \deleted{\unsure{on}} Pour une dim donnée $-1\le d\le n+1$, il y a \deleted{\unsure{au}} \inserted{un} \deleted{\unsure{seul}} seul des $Z_i,Z'_i$ de cette dim si $d=-1$ ou $d=n+1$ (savoir $Z_n^{\deleted{\prime}}$ et $Z'_{-1}$), et \inserted{exactement} deux si $0\le d\le n$, savoir $Z'_{n-d}=H_0\cap\cdots\cap H_{n-d}$ \deleted{\unsure{de}} et $Z_{n-d-1}^{\deleted{\prime}}=\mathrm{Lin}(h_n,\ldots,h_{n-d})$, \deleted{codim} et ils \deleted{\unsure{ne}} sont égaux sss $h_{n-d},h_{n-d+1}\in Z'_{n-d}$, i.e. $h_{n-d}\in H_{n-d}\cap H_{n-d-1}$ $h_{n-d+1}\in H_{n-d}$

% L13 -- lower left, framed note
\margin{\uline{NB} Si $c\notin C$ i.e. $\rho\ne0$, alors les $Z_i,Z'_i$ tous distincts, car $Z'_i\not\subset Z_j^{\deleted{\prime}}$ pour tts $i,j$}
% archive sheet number: 79
